\documentclass[11pt,letterpaper]{article}

\usepackage[spanish,es-noshorthands]{babel}
\usepackage{fontspec}
\setmainfont{Source Serif Pro}
\usepackage{microtype}
\usepackage{geometry}
\usepackage{booktabs}
\usepackage{tabularx}
\usepackage{enumitem}
\usepackage{xcolor}
\usepackage{graphicx}
\usepackage{csquotes}
\usepackage{amsmath}
\usepackage{siunitx}
\usepackage{tikz}
\usetikzlibrary{arrows.meta,positioning,shapes.geometric,fit,backgrounds}
\usepackage{xurl}
\usepackage[hidelinks]{hyperref}
\usepackage[backend=biber,style=apa,sorting=nyt]{biblatex}
\DeclareLanguageMapping{spanish}{spanish-apa}

\geometry{margin=2.5cm}
\setlength{\emergencystretch}{1em}
\setlength{\parindent}{0pt}
\setlength{\parskip}{0.65em}
\setlist{nosep}
\addbibresource{referencias.bib}
\AtBeginBibliography{\sloppy}

\sisetup{
  output-decimal-marker = {,},
  group-separator = {.},
  group-minimum-digits = 4,
  per-mode = symbol,
  detect-all
}
\DeclareSIUnit{\ton}{ton}
\DeclareSIUnit{\kWh}{kWh}
\DeclareSIUnit{\lumen}{lm}
\DeclareSIUnit{\bit}{bit}

\title{\textbf{La paradoja entre la curva S y la ley de Moore}\\
\large Por qué una tecnología puede madurar y progresar exponencialmente al
mismo tiempo}
\author{Carlos Luis Rojas Aragonés\\
Gestión del Desarrollo Tecnológico}
\date{29 de julio de 2026}

\begin{document}

\maketitle

\section{La paradoja}

Los dos modelos que la gestión tecnológica usa con más frecuencia parecen
contradecirse.

La \textbf{curva S} sostiene que el desempeño de una tecnología, medido en su
figura de mérito relevante, crece lentamente al principio, se acelera cuando se
resuelven los cuellos de botella principales y luego se aplana al aproximarse a
un límite físico o económico. Su implicación de gestión es inequívoca: toda
tecnología tiene un techo, la inversión adicional produce rendimientos
decrecientes en el tramo superior, y el momento de saltar a la siguiente curva
llega antes de que la curva actual haya dejado de mejorar
\parencite{foster1986}.

La \textbf{ley de Moore} afirma lo aparentemente opuesto. Formulada en 1965 sobre
apenas cinco puntos de datos, observó que el número de componentes en el circuito
integrado de costo unitario mínimo se duplicaba aproximadamente cada año, y
proyectó que la tendencia continuaría una década \parencite{moore1965}. En 1975
Moore revisó el ritmo a una duplicación cada dos años \parencite{moore1975}. La
proyección no duró diez años: duró más de cincuenta. Su implicación de gestión
también es inequívoca, y es la contraria: no planifique para un techo, planifique
para que el desempeño de mañana sea el doble del de hoy.

La contradicción es visible en la práctica: en 1971 el Intel 4004 integraba unos
\num{2300} transistores; en 2023 el Apple M2 Ultra integraba
\num{134000000000}. Son \num{25.8} duplicaciones en \num{52} años, es decir una
duplicación cada \num{2.02} años, exactamente el ritmo anunciado por Moore en
1975. Ninguna curva S ordinaria se comporta así durante medio siglo.

\begin{quote}
\textbf{Tesis de este trabajo.} La paradoja es aparente, no real. Surge de
comparar dos modelos que hablan de objetos distintos, con variables
independientes distintas y en escalas distintas. Correctamente alineados, la ley
de Moore no contradice la curva S: es \emph{el resultado que produce la curva S}
cuando el esfuerzo acumulado en I+D crece de forma exponencial y cuando cada
paradigma agotado es reemplazado por otro antes de saturarse. Y una vez
establecido eso, la ley de Moore deja de ser un caso excepcional: es un miembro
más de una familia amplia y bien documentada de regularidades exponenciales que
incluye desde la fabricación de fuselajes en 1936 hasta el precio de la luz
artificial desde el siglo XIV.
\end{quote}

\section{Qué afirma exactamente cada modelo}

Antes de resolver la contradicción conviene enunciar con precisión qué dice cada
modelo, porque la mayor parte del conflicto desaparece en el enunciado.

\begin{table}[ht]
  \centering
  \small
  \caption{Los dos modelos comparados en sus propios términos.}
  \label{tab:modelos}
  \begin{tabularx}{\textwidth}{@{}l >{\raggedright\arraybackslash}X >{\raggedright\arraybackslash}X@{}}
    \toprule
    & \textbf{Curva S (Foster)} & \textbf{Ley de Moore} \\
    \midrule
    Objeto que describe
      & Una tecnología concreta, un paradigma o un diseño dominante
      & Una industria completa y su trayectoria acumulada de paradigmas \\
    Variable dependiente
      & Una figura de mérito técnica
      & Densidad de transistores y las métricas derivadas de ella \\
    Variable independiente
      & \textbf{Esfuerzo acumulado} de ingeniería o inversión
      & \textbf{Tiempo} calendario \\
    Mecanismo
      & Rendimientos decrecientes al aproximarse a un límite físico
      & Escalado dimensional, aprendizaje por producción e inversión creciente \\
    Forma funcional
      & Logística o sigmoide
      & Exponencial, es decir recta en escala semilogarítmica \\
    Estatus epistemológico
      & Modelo con mecanismo causal explícito
      & Regularidad empírica, además con componente de profecía autocumplida \\
    Qué predice
      & Cuándo dejará de rendir la inversión y cuándo hay que cambiar de curva
      & El ritmo de mejora esperable en el próximo ciclo de producto \\
    \bottomrule
  \end{tabularx}
\end{table}

Dos observaciones de la Tabla~\ref{tab:modelos} bastan para desactivar buena
parte del problema.

La primera es que \textbf{la curva S de Foster no está dibujada contra el
tiempo}. Su eje horizontal es el esfuerzo acumulado de ingeniería, y Foster
insiste en ese punto justamente porque el tiempo es un sustituto engañoso: dos
tecnologías con el mismo calendario pueden estar en puntos muy diferentes de su
curva si una recibió cien veces más inversión que la otra \parencite{foster1986}.
La ley de Moore, en cambio, está formulada contra el tiempo. Comparar ambas
directamente equivale a comparar dos gráficos con abscisas incompatibles.

La segunda es que \textbf{no describen el mismo objeto}. La curva S describe un
paradigma; la ley de Moore describe una industria que ha cambiado de paradigma
al menos media docena de veces sin cambiar de figura de mérito. Son dos niveles
de análisis distintos, y el error de la paradoja es tratarlos como si fueran uno.

\section{Resolución: cuatro desajustes que hay que corregir}

\subsection*{01. Desajuste de nivel: la envolvente no es una curva S, es una
sucesión de curvas S}

La ley de Moore no describe el ciclo de vida del transistor bipolar, ni el del
NMOS, ni el del CMOS planar, ni el del FinFET. Describe la \emph{envolvente} de
todos ellos. Cada uno de esos paradigmas sí siguió una curva S propia: entrada
lenta, tramo de mejora rápida, saturación y sustitución. Lo que la industria de
semiconductores logró, y lo que la hace notable, no fue evitar la saturación,
sino \textbf{encadenar sustituciones con la puntualidad suficiente para que la
envolvente no tuviera discontinuidades visibles}.

Foster ya describía este fenómeno con el nombre de \emph{curvas en cascada}, y
Sahal lo formalizó como \enquote{guías tecnológicas}: familias de curvas S
sucesivas que comparten una misma trayectoria de mejora y una misma figura de
mérito \parencite{foster1986,sahal1985}. Utterback documentó el mismo patrón en
la sucesión de diseños dominantes de varias industrias \parencite{utterback1994}.

\begin{figure}[ht]
  \centering
  \begin{tikzpicture}[xscale=1.28,yscale=0.92]
    % ejes
    \draw[-{Stealth[length=2mm]}] (-0.2,-0.2) -- (9.9,-0.2)
      node[below left=1mm and 0mm] {\footnotesize Tiempo};
    \draw[-{Stealth[length=2mm]}] (-0.2,-0.2) -- (-0.2,6.6)
      node[above, rotate=0, align=center]
      {\footnotesize $\log$ (figura de mérito)};
    % envolvente
    \draw[dashed, thick, red!65!black, domain=0.1:9.3, samples=2, variable=\x]
      plot (\x, {0.75*\x - 0.125})
      node[right, align=left, text width=1.6cm]
      {\footnotesize\color{red!65!black} Ley de\\ Moore};
    % curvas S
    \foreach \i/\name in {0/P1, 1/P2, 2/P3, 3/P4} {
      \pgfmathsetmacro{\xc}{1.5 + 2*\i}
      \pgfmathsetmacro{\yc}{1.0 + 1.5*\i}
      \draw[thick, blue!60!black, domain=\xc-1.7:\xc+1.7, samples=60,
        variable=\x, smooth]
        plot (\x, {\yc - 1.0 + 2.0/(1+exp(-2.2*(\x-\xc)))});
      \node[blue!60!black, font=\scriptsize] at (\xc-1.25, \yc+1.15) {\name};
    }
  \end{tikzpicture}
  \caption{La ley de Moore como envolvente de curvas S en cascada. Cada
  paradigma tecnológico ($P1$ a $P4$) satura, pero su sucesor entra antes de que
  la saturación se haga visible en la agregación. La recta en escala
  semilogarítmica no describe ninguna tecnología concreta: describe la política
  de reemplazo de una industria. Elaboración propia.}
  \label{fig:envolvente}
\end{figure}

La consecuencia de gestión es importante y con frecuencia se pasa por alto: la
envolvente \textbf{no es un mecanismo, es un resultado agregado}. No existe
ninguna ley física que garantice que un paradigma sucesor estará disponible
cuando el actual se agote. La envolvente exponencial es un hecho retrospectivo
sobre una industria que invirtió lo necesario para que hubiera sucesor. Este es
el punto donde la lectura de Kurzweil, que eleva la cascada a \enquote{ley de
retornos acelerados} con validez general \parencite{kurzweil2005}, va más allá de
lo que la evidencia sostiene: extrapolar la envolvente supone que la cascada
continúa, y esa es una hipótesis sobre inversión y sobre la existencia de un
sucesor viable, no una consecuencia de los datos históricos.

\subsection*{02. Desajuste de escala: la exponencial es el tramo de despegue de
la propia curva S}

El segundo desajuste es puramente geométrico y explica por qué el conflicto
resulta tan persuasivo a la vista. Sea la logística estándar con techo $L$:

\begin{equation}
  y(t) = \frac{L}{1 + e^{-r(t - t_0)}}
  \label{eq:logistica}
\end{equation}

Cuando $y \ll L$, el denominador está dominado por el término exponencial y la
expresión se reduce a

\begin{equation}
  y(t) \approx L\, e^{\,r(t-t_0)}
  \label{eq:aprox}
\end{equation}

Es decir: \textbf{el tramo de despegue de una curva S es indistinguible de una
exponencial pura}. En escala semilogarítmica, que es la escala en la que siempre
se dibuja la ley de Moore, la logística aparece como una recta durante toda la
fase en que el desempeño se mantiene lejos del techo, y solo se dobla cuando ya
está cerca de él. Una tecnología puede pasar décadas en ese tramo recto y ser,
sin ninguna contradicción, una curva S.

\begin{figure}[ht]
  \centering
  \begin{tikzpicture}[yscale=0.62]
    \begin{scope}
      \draw[-{Stealth[length=2mm]}] (0,0) -- (5.3,0)
        node[below left=1mm and 0mm] {\footnotesize $t$};
      \draw[-{Stealth[length=2mm]}] (0,0) -- (0,7.1)
        node[above] {\footnotesize $y$};
      \draw[dotted, gray] (0,6) -- (5.1,6) node[right, font=\scriptsize] {$L$};
      \draw[thick, blue!60!black, domain=0:5, samples=80, variable=\x, smooth]
        plot (\x, {6/(1+exp(-1.8*(\x-2.5)))});
      \node[font=\scriptsize, align=center] at (2.5,-1.0)
        {Escala lineal:\\ la S es evidente};
    \end{scope}
    \begin{scope}[xshift=7.4cm]
      \draw[-{Stealth[length=2mm]}] (0,0) -- (5.3,0)
        node[below left=1mm and 0mm] {\footnotesize $t$};
      \draw[-{Stealth[length=2mm]}] (0,0) -- (0,7.1)
        node[above] {\footnotesize $\log y$};
      \draw[dotted, gray] (0,6) -- (5.1,6) node[right, font=\scriptsize] {$\log L$};
      \draw[thick, blue!60!black, domain=0:5, samples=80, variable=\x, smooth]
        plot (\x, {6 - ln(1+exp(-1.8*(\x-2.5)))*0.98});
      \draw[dashed, red!65!black] (0.15,0.42) -- (2.9,5.3);
      \node[font=\scriptsize, align=center] at (2.5,-1.0)
        {Escala semilogarítmica:\\ recta hasta muy cerca del techo};
    \end{scope}
  \end{tikzpicture}
  \caption{La misma curva logística en dos escalas. La recta discontinua marca
  el tramo en el que la curva S es empíricamente indistinguible de una
  exponencial. Una serie observada dentro de ese tramo admite las dos lecturas.
  Elaboración propia.}
  \label{fig:escalas}
\end{figure}

Esto tiene una consecuencia metodológica directa: \textbf{la observación de una
recta en semilogarítmico no permite descartar la existencia de un techo}. Solo
permite afirmar que el techo todavía no se ha alcanzado. Determinar la distancia
al techo exige información externa a la serie, típicamente un cálculo del límite
físico, como el que se aplica al horno de arco eléctrico cuando se compara su
consumo real con el mínimo termodinámico de fusión del acero. Sin ese cálculo,
extrapolar la recta es una apuesta, no una predicción.

\subsection*{03. Desajuste de variable: la ley de Moore es la ley de Wright más
inversión exponencial}

Este es el núcleo analítico de la resolución, y proviene de un resultado que
Sahal publicó en 1979 \parencite{sahal1979} y que Nagy y colaboradores
reformularon y contrastaron empíricamente décadas después
\parencite{nagy2013,farmer2016}.

La \textbf{ley de Wright}, formulada en 1936 a partir de datos de fabricación de
fuselajes, establece que el costo unitario cae en un porcentaje fijo cada vez que
se duplica la producción acumulada \parencite{wright1936}. En forma de potencia:

\begin{equation}
  y(x) = y_0 \left(\frac{x}{x_0}\right)^{-w}
  \label{eq:wright}
\end{equation}

donde $x$ es el esfuerzo o la producción acumulada, $y$ la figura de mérito
(costo por unidad, energía por operación, precio por watt) y $w > 0$ la
elasticidad de aprendizaje. Nótese que la variable independiente de la
ecuación~\eqref{eq:wright} es la misma que la de la curva S de Foster: esfuerzo
acumulado, no tiempo.

Supongamos ahora que el esfuerzo acumulado crece exponencialmente en el tiempo,
lo que en semiconductores es una descripción razonable de lo que ocurrió, dada
la expansión del mercado, del gasto en I+D y del capital por planta:

\begin{equation}
  x(t) = x_0\, e^{\,g t}
  \label{eq:esfuerzo}
\end{equation}

Sustituyendo \eqref{eq:esfuerzo} en \eqref{eq:wright}:

\begin{equation}
  \boxed{\;y(t) = y_0\, e^{-w g t}\;}
  \label{eq:moore}
\end{equation}

La ecuación~\eqref{eq:moore} \textbf{es} la ley de Moore. Su ritmo de mejora es
el producto $m = w \cdot g$ de dos factores de naturaleza completamente
distinta: una elasticidad técnica de aprendizaje y una tasa económica de
crecimiento del esfuerzo. De aquí se siguen tres conclusiones que resuelven la
paradoja sin necesidad de abandonar ninguno de los dos modelos:

\begin{enumerate}
  \item \textbf{La exponencialidad en el tiempo no es una propiedad de la
    tecnología.} Es el producto de una propiedad de la tecnología ($w$) y una
    decisión de la industria ($g$). Una tecnología con excelente aprendizaje y
    esfuerzo estancado no exhibirá ninguna ley de Moore.
  \item \textbf{La curva S y la ley de Moore son el mismo modelo con distintos
    supuestos sobre $x(t)$.} Si el esfuerzo acumulado crece exponencialmente, se
    observa una exponencial en el tiempo. Si el esfuerzo acumulado sigue una
    difusión logística, que es lo habitual en tecnologías con mercados que
    saturan, se observa una curva S en el tiempo aunque el aprendizaje $w$ sea
    constante. La forma temporal la impone la trayectoria de la inversión, no la
    física del dispositivo.
  \item \textbf{Hay tres formas independientes de matar una exponencial.} Que
    $w \to 0$, porque la figura de mérito se acerca a su límite físico y el
    aprendizaje se agota. Que $g \to 0$, porque el mercado o el capital dejan de
    crecer exponencialmente. O que el paradigma actual se agote sin sucesor
    disponible, rompiendo la cascada de la Figura~\ref{fig:envolvente}. Las tres
    son visibles hoy en semiconductores, y la segunda es la menos discutida: el
    costo de una planta de nodo avanzado superó los veinte mil millones de
    dólares, contra unos pocos cientos de millones a comienzos de los años
    noventa, y el número de empresas capaces de fabricar en el nodo más avanzado
    cayó de más de veinte a tres. Moore mismo anticipó esa restricción económica
    \parencite{moore1995,waldrop2016}.
\end{enumerate}

\subsection*{04. Desajuste de figura de mérito: la curva S ya llegó, solo que no
a la métrica que se cita}

El cuarto desajuste es empírico y es, en mi opinión, el más contundente. La
premisa de que la ley de Moore \enquote{sigue vigente} solo se sostiene si se
elige una figura de mérito muy específica: transistores por chip. En cuanto se
examinan las otras métricas de la misma tecnología, la saturación predicha por la
curva S aparece con puntualidad, en algunos casos hace veinte años.

\begin{table}[ht]
  \centering
  \small
  \caption{Saturación diferencial por figura de mérito en semiconductores. La
  curva S no falló: llegó a distinto tiempo a cada métrica.}
  \label{tab:saturacion}
  \begin{tabularx}{\textwidth}{@{}>{\raggedright\arraybackslash}X l l l@{}}
    \toprule
    \textbf{Figura de mérito} & \textbf{Ritmo histórico} &
    \textbf{Ritmo reciente} & \textbf{Estado} \\
    \midrule
    Transistores por chip
      & $\times 2$ / 2 años & $\times 2$ / 3 años aprox. & Activo, desacelerando \\
    Escalado de Dennard (densidad de potencia constante)
      & Vigente 1974--2004 & Roto & Terminado \\
    Frecuencia de reloj
      & $\times 2$ / 2,5 años & Plana en 3--5 GHz & Saturado desde 2004 \\
    Rendimiento de un solo hilo
      & \SI{52}{\percent}/año (1986--2003) & \SI{3.5}{\percent}/año (2015--2018)
      & Saturado \\
    Costo por transistor
      & Descenso sostenido & Plano o creciente & Invertido \\
    Costo de capital por planta
      & Creciente & $>$ 20.000 M USD & Límite económico \\
    Densidad areal en disco magnético
      & \SI{40}{\percent}/año (años 1990) & $<$ \SI{10}{\percent}/año & Roto \\
    Rendimiento por watt en aceleradores especializados
      & Curva nueva & Alto & Activo \\
    \bottomrule
  \end{tabularx}
\end{table}

Los datos de la Tabla~\ref{tab:saturacion} son elocuentes. El escalado de
Dennard, que era el mecanismo físico que convertía la miniaturización en
velocidad sin aumentar la densidad de potencia, dejó de funcionar alrededor de
2005 al volverse dominantes las corrientes de fuga \parencite{dennard1974}. La
frecuencia de reloj se detuvo en el mismo momento. El rendimiento de un solo
hilo, que creció al \SI{52}{\percent} anual entre 1986 y 2003, cayó al
\SI{23}{\percent} entre 2003 y 2011, al \SI{12}{\percent} entre 2011 y 2015 y al
\SI{3.5}{\percent} entre 2015 y 2018 \parencite{hennessy2019}. Ese perfil, un
factor de quince en la tasa de mejora en tres décadas, es exactamente el tramo
superior de una curva S.

Lo que ocurrió después es, además, la respuesta canónica de la teoría de la curva
S: la industria \textbf{cambió de curva}. Abandonó la mejora del procesador de
propósito general y se movió a la especialización de la arquitectura, a los
aceleradores de dominio específico, al apilamiento tridimensional, a los
\emph{chiplets} y a la optimización del software
\parencite{hennessy2019,leiserson2020,thompson2021}. La aparente continuidad de
la ley de Moore en los últimos quince años no es la prolongación de una curva:
es la superposición de una nueva.

\section{La ley de Moore no es un caso especial}

La segunda parte del planteamiento pide identificar otra ley o teoría que
muestre que la ley de Moore no es excepcional. La respuesta corta es que la ley
de Moore es \textbf{un caso particular de la ley de Wright}, que es anterior en
veintinueve años, más general y con mejor desempeño predictivo. La respuesta
larga incluye la evidencia sistemática acumulada sobre docenas de dominios
tecnológicos.

\subsection*{01. La ley de Wright y la evidencia sistemática}

Wright publicó su curva de aprendizaje en 1936, estudiando el costo de
fabricación de fuselajes, y encontró una reducción cercana al
\SI{20}{\percent} del costo unitario por cada duplicación de la producción
acumulada \parencite{wright1936}. No es una ley sobre electrónica: es una ley
sobre el aprendizaje organizacional y tecnológico en general.

La comparación directa entre ambas formulaciones ha sido hecha con rigor
estadístico. Nagy y colaboradores contrastaron las dos hipótesis sobre una base
de \num{62} tecnologías y encontraron que ambas se ajustan bien, con una ligera
ventaja para Wright, y que las dos son consistentes entre sí precisamente por la
identidad de la ecuación~\eqref{eq:moore} \parencite{nagy2013}. Farmer y Lafond
formalizaron después los intervalos de predicción sobre \num{53} tecnologías y
mostraron que el error de pronóstico crece con el horizonte de una manera
estimable, lo que convierte estas curvas en herramientas de planificación con
incertidumbre cuantificada y no en profecías \parencite{farmer2016}. Magee y
colaboradores midieron las tasas anuales de mejora en \num{28} dominios
tecnológicos y encontraron mejora exponencial sostenida en todos ellos, con
ritmos que van de unos pocos puntos porcentuales anuales hasta bastante más del
\SI{30}{\percent}, siendo las tecnologías de información las más rápidas pero no
las únicas \parencite{magee2016,koh2006}.

El hallazgo agregado de esa literatura es el que interesa aquí: \textbf{la mejora
exponencial sostenida es el comportamiento típico de una tecnología en su tramo
de despegue, no la excepción}. Lo excepcional de los semiconductores es la
magnitud del ritmo y su duración, no su forma funcional.

\subsection*{02. Un catálogo de exponenciales}

La Tabla~\ref{tab:leyes} recoge regularidades del mismo tipo en dominios muy
distintos. Se incluye deliberadamente tanto casos vigentes como casos
interrumpidos, porque los segundos son los que demuestran que la curva S termina
imponiéndose.

\begin{table}[ht]
  \centering
  \small
  \caption{Regularidades exponenciales documentadas en distintos dominios
  tecnológicos.}
  \label{tab:leyes}
  \begin{tabularx}{\textwidth}{@{}>{\raggedright\arraybackslash}p{0.19\textwidth} >{\raggedright\arraybackslash}X >{\raggedright\arraybackslash}p{0.20\textwidth} l@{}}
    \toprule
    \textbf{Regularidad} & \textbf{Figura de mérito y ritmo} &
    \textbf{Dominio} & \textbf{Estado} \\
    \midrule
    Ley de Wright (1936)
      & Costo unitario, $-$\SI{20}{\percent} por duplicación acumulada
      & Fuselajes de aeronaves & Vigente como modelo general \\
    Ley de Moore (1965)
      & Transistores por chip, $\times 2$ cada 2 años
      & Circuitos integrados & Desacelerando \\
    Ley de Koomey
      & Cómputo por unidad de energía, $\times 2$ cada 1,57 años (1946--2009)
      & Eficiencia energética del cómputo & Desacelerando \\
    Ley de Kryder
      & Densidad areal, cerca del \SI{40}{\percent} anual en los años 1990
      & Almacenamiento magnético & Roto tras 2010 \\
    Ley de Butter
      & Costo por bit transmitido, mitad cada 9 meses aprox.
      & Fibra óptica & Vigente \\
    Ley de Haitz
      & Costo por lumen $/10$ y flujo por dispositivo $\times 20$ por década
      & Iluminación de estado sólido & Vigente \\
    Curva de Carlson
      & Costo por genoma secuenciado, más rápido que Moore desde 2007
      & Secuenciación de ADN & Desacelerando \\
    Ley de Swanson
      & Precio por watt, cerca de $-$\SI{20}{\percent} por duplicación
      & Módulos fotovoltaicos & Vigente \\
    Aprendizaje en baterías
      & Precio por kWh, $-$\SI{20}{\percent} a $-$\SI{25}{\percent} por duplicación
      & Celdas de ion litio & Vigente \\
    Precio de la luz
      & Precio real por lumen-hora, descenso sostenido durante siete siglos
      & Servicios de iluminación & Vigente \\
    Velocidad de crucero
      & Estancada cerca de Mach 0,85 desde 1958
      & Aviación comercial & Saturado \\
    \bottomrule
  \end{tabularx}
\end{table}

Cuatro de estas entradas merecen comentario porque cada una aporta un argumento
distinto contra la excepcionalidad de la ley de Moore.

\textbf{La ley de Koomey} \parencite{koomey2011} muestra que el cómputo por
unidad de energía se duplicó cada \num{1.57} años entre 1946 y 2009, un intervalo
que abarca válvulas, transistores discretos y circuitos integrados. Es decir: la
regularidad es anterior al objeto que la ley de Moore describe y sobrevive al
cambio de paradigma físico. Eso es una envolvente en el sentido de la
Figura~\ref{fig:envolvente}, y es una evidencia fuerte de que lo que se está
midiendo es una trayectoria funcional y no una tecnología.

\textbf{La ley de Kryder} es el contraejemplo decisivo. La densidad areal del
disco magnético creció a un ritmo comparable o superior al de la ley de Moore
durante los años noventa, hasta que las limitaciones del límite superparamagnético
la frenaron por debajo del \SI{10}{\percent} anual en la década de 2010
\parencite{fontana2018}. Una exponencial de la misma familia, en la misma
industria y en el mismo período, terminó. Quien sostenga que la ley de Moore es
inmune al techo de la curva S debe explicar por qué la ley de Kryder no lo fue.

\textbf{El precio de la luz artificial} es el ejemplo histórico más largo y, para
mi propósito, el más ilustrativo. Nordhaus reconstruyó el precio real del
servicio de iluminación desde la antigüedad y mostró un descenso de varios
órdenes de magnitud, con una caída del orden de mil veces solo entre 1800 y 1992
\parencite{nordhaus1996}. Fouquet y Pearson extendieron el ejercicio al Reino
Unido entre 1300 y 2000 y encontraron un descenso del orden de tres mil veces en
el precio real por lumen-hora \parencite{fouquet2006}. Esa trayectoria de siete
siglos es exponencial en su conjunto y está compuesta por una sucesión limpia de
curvas S: vela de sebo, vela de cera, aceite de ballena, gas de hulla, queroseno,
arco eléctrico, incandescencia, fluorescencia y diodo emisor de luz. Cada una
saturó; la envolvente no. Es la ley de Moore, seiscientos años antes y sin
silicio. Y su capítulo actual, el LED, tiene su propia formulación exponencial en
la ley de Haitz \parencite{haitz2011}.

\textbf{La aviación comercial} cierra el argumento por el lado opuesto. La
velocidad de crucero de los aviones de pasajeros creció de forma acelerada
durante medio siglo y se detuvo cerca de Mach 0,85 con la generación del Boeing
707, a finales de los años cincuenta, donde permanece. El límite no fue una
imposibilidad física, dado que el Concorde voló, sino la combinación de la
barrera transónica de arrastre con la economía del combustible. La industria
cambió de figura de mérito, de velocidad a costo por asiento-kilómetro, que es
exactamente lo que hicieron los semiconductores en 2005 al pasar de frecuencia de
reloj a rendimiento por watt.

\subsection*{03. Un elemento adicional: el componente social de la ley de Moore}

Hay una razón, ajena a la física y a la economía, por la que la ley de Moore
parece más robusta de lo que es: durante décadas funcionó como \textbf{mecanismo
de coordinación de una industria entera}. La hoja de ruta sectorial, primero la
ITRS y después la IRDS, convirtió la proyección en un objetivo compartido por
fabricantes de equipos, de materiales, de herramientas de diseño y de chips
\parencite{mack2011,irds2022}. Cuando miles de organizaciones planifican su
inversión para cumplir una curva, la curva se cumple. Eso hace de la ley de
Moore, en parte, una profecía autocumplida, y explica por qué el término $g$ de
la ecuación~\eqref{eq:moore} se mantuvo alto tanto tiempo. También explica por
qué su desaceleración coincide con el momento en que la coordinación se volvió
demasiado costosa para casi todos los participantes.

\section{Dónde está el techo de la ley de Moore}

Si la ley de Moore es una curva S en su tramo recto, cabe preguntar dónde está su
asíntota. La respuesta se obtiene con el mismo método que se emplea para calcular
el mínimo termodinámico de un proceso industrial: identificar la barrera física y
medir la distancia que separa el desempeño actual de ella.

Para la lógica digital irreversible, el límite fundamental es el de Landauer:
borrar un bit de información disipa como mínimo una energía
$E_{min} = k_B T \ln 2$ \parencite{landauer1961}. A temperatura ambiente
($T = \SI{300}{\kelvin}$):

\begin{equation}
  E_{min} = \SI{1.38e-23}{\joule\per\kelvin} \times \SI{300}{\kelvin}
  \times \ln 2 \approx \SI{2.9e-21}{\joule}
  \label{eq:landauer}
\end{equation}

Un conmutador CMOS actual disipa del orden de \SI{1e-17}{\joule} por transición,
lo que sitúa la tecnología entre tres y cuatro órdenes de magnitud por encima de
su límite fundamental \parencite{zhirnov2003}. La lectura correcta de ese
resultado es doble y aparentemente ambivalente, pero no lo es:

\begin{itemize}
  \item \textbf{El techo existe.} No es una conjetura de gestión, es
    termodinámica. La ley de Moore tiene una asíntota, igual que el consumo
    eléctrico del horno de arco eléctrico tiene un mínimo de fusión.
  \item \textbf{El techo está lejos en términos de física y cerca en términos de
    ingeniería.} Quedan tres o cuatro órdenes de magnitud teóricos, pero el
    escalado dimensional que los entregaba ya se agotó, y las barreras efectivas
    que detuvieron el progreso desde 2005 fueron la densidad de potencia, la
    litografía y el costo de capital, no el límite de Landauer. Es la distinción
    clásica entre el límite físico y el límite económicamente accesible.
\end{itemize}

Esta es, en el fondo, la misma lección que ofrece el cálculo del mínimo
termodinámico de cualquier proceso: el límite teórico define el espacio de lo
posible, pero el límite relevante para la decisión de inversión es aquel a partir
del cual el costo marginal de la siguiente unidad de mejora crece más rápido que
su valor.

\section{Síntesis: cómo resolvería la paradoja}

Mi resolución, formulada como marco operativo, tiene cinco reglas.

\begin{enumerate}
  \item \textbf{Distinguir el nivel de análisis antes de dibujar la curva.} Una
    curva S describe un paradigma; una ley exponencial describe una envolvente.
    Nunca deben trazarse en el mismo gráfico sin decir cuál es cuál. El error de
    la paradoja consiste en comparar el ciclo de vida de un componente con la
    trayectoria agregada de una industria.
  \item \textbf{Elegir la figura de mérito antes de opinar sobre madurez.} La
    Tabla~\ref{tab:saturacion} demuestra que la misma tecnología puede estar
    saturada en una métrica y en pleno despegue en otra. La pregunta
    \enquote{¿está madura esta tecnología?} está mal planteada; la correcta es
    \enquote{¿qué figura de mérito está madura y cuál no?}.
  \item \textbf{Calcular el límite físico y medir la brecha.} Sin un cálculo del
    techo, la Figura~\ref{fig:escalas} muestra que no es posible distinguir un
    despegue de una recta indefinida. El límite de Landauer, el mínimo
    termodinámico de fusión o el límite de Shannon para un canal cumplen esa
    función: convierten una extrapolación en una medición de recorrido restante.
  \item \textbf{Descomponer el ritmo observado en aprendizaje e inversión.} La
    ecuación~\eqref{eq:moore} indica que $m = w \cdot g$. Un ritmo exponencial
    sostenido por $g$ alto y $w$ decreciente es frágil, porque depende de que la
    inversión siga creciendo. Distinguir los dos términos permite anticipar la
    inflexión antes de verla en la serie.
  \item \textbf{Vigilar la cascada, no la envolvente.} La envolvente solo
    continúa si existe un paradigma sucesor en desarrollo. La señal de alerta
    útil no es la desaceleración de la envolvente, que llega tarde, sino la
    ausencia de un sucesor viable en el momento en que el paradigma actual entra
    en su tramo superior. Es la formulación operativa de la advertencia de Foster
    y el mismo mecanismo que explica la disrupción en Christensen
    \parencite{foster1986,christensen1997}.
\end{enumerate}

\begin{figure}[ht]
  \centering
  \begin{tikzpicture}[
    node distance=4mm,
    caja/.style={draw, rounded corners, align=center, minimum height=13mm,
      text width=0.205\textwidth, font=\small},
    flecha/.style={-{Stealth[length=2.2mm]}, thick}
  ]
    \node[caja, fill=blue!6] (nivel) {¿Paradigma\\ o envolvente?};
    \node[caja, right=of nivel, fill=blue!8] (fom) {¿Qué figura\\ de mérito?};
    \node[caja, right=of fom, fill=orange!12] (limite) {¿Cuál es el\\ límite físico?};
    \node[caja, right=of limite, fill=green!10] (decision) {$m = w \cdot g$:\\
      ¿optimizar o saltar?};
    \draw[flecha] (nivel) -- (fom);
    \draw[flecha] (fom) -- (limite);
    \draw[flecha] (limite) -- (decision);
  \end{tikzpicture}
  \caption{Secuencia de decisión para resolver la paradoja en un caso concreto.
  Elaboración propia.}
  \label{fig:marco}
\end{figure}

\section{Limitaciones del argumento}

Cuatro precisiones honestas sobre el alcance de lo expuesto.

\begin{enumerate}
  \item \textbf{Sesgo de supervivencia en la evidencia.} Las bases de datos de
    curvas de progreso recogen tecnologías que prosperaron lo suficiente para ser
    medidas durante décadas \parencite{nagy2013,farmer2016}. Las tecnologías que
    se estancaron temprano están subrepresentadas, de modo que la afirmación
    \enquote{la mejora exponencial es común} está probablemente sobreestimada en
    magnitud, aunque no en su punto central: la ley de Moore no es única.
  \item \textbf{La identidad de la ecuación~\eqref{eq:moore} es una explicación
    parcial.} Requiere conocer $x(t)$, y el esfuerzo acumulado en I+D es difícil
    de medir de forma consistente entre industrias. La derivación aclara la
    lógica de la relación mejor de lo que permite calibrarla numéricamente.
  \item \textbf{La elección de la figura de mérito condiciona la conclusión.}
    Este trabajo la usa como instrumento analítico, pero también podría usarse
    para fabricar la respuesta deseada. Un cambio de métrica puede prolongar
    artificialmente una envolvente que, para el usuario final, ya se detuvo. La
    disciplina consiste en fijar la métrica que importa al cliente antes de
    evaluar la madurez, no después.
  \item \textbf{El futuro de la cascada es una hipótesis, no un dato.} La
    computación cuántica, la fotónica integrada, los materiales bidimensionales o
    la computación reversible podrían constituir el próximo paradigma de la
    envolvente. Ninguno está hoy en condiciones de sostenerla, y afirmar lo
    contrario sería repetir el error de extrapolar la envolvente que se critica
    en este documento.
\end{enumerate}

\section*{Conclusiones}

La paradoja entre la curva S y la ley de Moore se disuelve al advertir que los
dos modelos no compiten por describir el mismo objeto. La curva S describe un
paradigma tecnológico en función del esfuerzo acumulado; la ley de Moore describe
la envolvente de una sucesión de paradigmas en función del tiempo.

La relación entre ambas es analítica y no meramente conceptual. Combinando la ley
de Wright con un crecimiento exponencial del esfuerzo acumulado se obtiene
exactamente una exponencial en el tiempo, con ritmo $m = w \cdot g$. La ley de
Moore es, por tanto, un caso particular de la curva de aprendizaje bajo un
supuesto específico sobre la trayectoria de la inversión. Cuando ese supuesto
falla, y $g$ deja de crecer, la exponencial temporal se convierte en una curva S
sin que la física del dispositivo haya cambiado.

La evidencia de que la ley de Moore no es especial es abundante. La ley de Wright
la antecede en veintinueve años y la generaliza. Los estudios sistemáticos sobre
decenas de dominios muestran mejora exponencial como comportamiento típico del
tramo de despegue. La ley de Koomey documenta la misma trayectoria a través de
tres paradigmas físicos distintos. La ley de Kryder recorrió una trayectoria
equivalente y se detuvo. Y el precio de la luz artificial describe, durante siete
siglos y sin electrónica alguna, la misma envolvente exponencial construida sobre
una cascada de curvas S sucesivas.

Finalmente, la curva S ya se manifestó en los semiconductores: en el escalado de
Dennard hacia 2005, en la frecuencia de reloj, en el rendimiento de un solo hilo
que pasó del \SI{52}{\percent} al \SI{3.5}{\percent} anual, y en el costo por
transistor. Lo que se mantuvo fue el conteo de transistores, sostenido primero
por cambios de arquitectura del dispositivo y luego por la especialización de la
computación, es decir, por curvas nuevas. Para la gestión tecnológica la
conclusión es que ningún modelo debe aplicarse sin declarar antes el nivel de
análisis, la figura de mérito y el límite físico. Hecho eso, la curva S y la ley
de Moore dejan de contradecirse y pasan a ser dos vistas del mismo fenómeno: una
mira el componente, la otra mira la cadena de reemplazos que lo sustituye.

\printbibliography[title={Referencias bibliográficas}]

\end{document}
